\section{Eigenvalue Lens: comprimir una enorme vocabulary matrix en una copying signature}

La sección anterior ya plegó cada head en token-to-token matrices, pero eso solo volvió el problema definible, todavía no legible: el lado de $M_{OV}$ y de $M_{QK}$ es igual al vocabulary size, por lo que enumerar entrada por entrada es costoso y además difícilmente permite un juicio global. Esta sección introduce el eigenvalue lens; la pregunta central es si con unas pocas spectral signatures se puede determinar si un head se parece más a copying, a anti-copying o a un different-token mapping. Al leer las figuras hay que recordar siempre que esta compresión pierde los tokens concretos, la basis y el contexto no lineal.

Si una matriz mapea un espacio vectorial de vuelta sobre sí mismo, pueden usarse eigenvectors/eigenvalues para resumir si en ciertas directions amplifica, invierte o rota.

\lecturefigure{slide-26-eigenvector-introduction.jpg}{Cuando la entrada y la salida pertenecen al mismo espacio, los eigenvectors pueden comprimir el comportamiento de la matriz.}{Video oficial de Stanford Online 00:33:58--00:34:21.}

\[
Mv=\lambda v.
\]

\begin{itemize}
  \item $v$: el mode cuya dirección no cambia tras aplicar la matriz;
  \item $\lambda$: el cambio de escala/fase de ese mode;
  \item $\lambda$ real positive: refuerzo en la misma dirección;
  \item $\lambda$ real negative: supresión en sentido inverso; $\lambda$ complex: rotación y escalamiento en un subespacio bidimensional.
\end{itemize}

\lecturefigure{slide-27-complex-eigenvalues.jpg}{Los token-to-token maps no simétricos pueden tener complex eigenvalues.}{Video oficial de Stanford Online 00:34:22--00:35:09.}

\subsection{OV spectrum y copying}

Para un OV map en token-space, los eigenmodes positivos indican que, al ver cierto grupo de token directions, se elevan los logits de output directions del mismo tipo, lo cual puede interpretarse como una copying tendency; los modes negativos son anti-copying; los modes complex/imaginary tienden más a mapear un tipo de token direction a otro distinto.

\lecturefigure{slide-28-ov-eigenvalue-meaning.jpg}{Explicación didáctica de la figura de OV eigenvalues: anti-copying, different-token mapping y copying.}{Video oficial de Stanford Online 00:35:10--00:35:55.}

\begin{warningbox}{El spectrum es un summary, no un diccionario semántico completo}
El embedding/unembedding del vocabulary no es necesariamente ortogonal, y la matriz también puede ser non-normal; los eigenvectors no necesariamente corresponden a un solo token nombrable. Reducir un eigenvalue positivo a copying es un diagnóstico práctico para small attention-only circuits, no un teorema semántico universal para cualquier matriz.
\end{warningbox}

\subsection{Distribución de copying en heads de una y dos capas}

Olah muestra que, en el modelo de una capa, el OV spectrum de diez (de doce) heads cae principalmente en la dirección real positiva, lo que indica que el modelo usa ampliamente la estrategia simple de «al ver un token, elevar el mismo token o uno relacionado». El modelo de dos capas también tiene muchos copying heads, pero el routing de varias capas permite que se activen con un pattern más preciso.

\lecturefigure{slide-29-layer-one-copying-heads.jpg}{Modelo de una capa: 10/12 attention heads tienen una copying signature fuerte.}{Video oficial de Stanford Online 00:35:56--00:36:05.}

\lecturefigure{slide-30-layer-two-copying-heads.jpg}{Modelo de dos capas: la segunda capa sigue teniendo varios copying heads fuertes.}{Video oficial de Stanford Online 00:36:06--00:36:41.}

\lecturefigure{slide-31-head-eigenvalue-histogram.jpg}{Compresión de la OV spectral signature en un histogram de heads de todo el modelo.}{Video oficial de Stanford Online 00:36:42--00:37:21.}

\teachervoice{El ponente enfatiza que esta es una bottom-up methodology: primero se deriva un summary legible a partir de la estructura de la matriz y después se afirma que «la mayoría de estos heads hacen copy», en lugar de suponer de antemano que todos los heads siguen cierto pattern y luego escoger ejemplos que lo respalden.}

\subsection{QK spectrum y límites de validez del método}

El positive mode de QK tiende a emparejar la misma token direction; el negative mode tiende a evitar el mismo token; el complex mode tiende a una relación entre tokens distintos. Esto tiene más sentido en modelos de varias capas, porque los previous heads pueden escribir en el residual actual información como «cuál fue el token anterior», y luego el QK de la capa siguiente realiza un same-feature matching.

\lecturefigure{slide-32-qk-eigenvalues-multilayer.jpg}{Los QK eigenvalues son más útiles en chains de varias capas: las capas anteriores mueven la información y las posteriores la emparejan.}{Video oficial de Stanford Online 00:37:22--00:37:37.}

\lecturefigure{slide-33-eigenanalysis-breaks-at-scale.jpg}{La naive eigenvalue analysis deja de funcionar al crecer el modelo y al agregar MLP.}{Video oficial de Stanford Online 00:37:38--00:38:35.}

\begin{warningbox}{El MLP invalida la idea de que «una matriz lo explica todo»}
Un attention-only model puede linealizarse localmente con un fixed pattern; el MLP introduce activation-dependent nonlinear feature transforms, y LayerNorm también acopla las coordenadas. En modelos grandes todavía se pueden analizar el Jacobian local, las features y los paths, pero no puede copiarse tal cual la small-model eigenvalue story.
\end{warningbox}

\teachervoice{Olah dice explícitamente que el naive spectral method «completely breaks down» al agregar MLP. Una nota de calidad debe conservar esta frase y no presentar las eigenvalue plots como una caja de herramientas general de interpretabilidad para LLM.}

\subsection{Resumen de la sección}

El eigenvalue lens comprime un enorme token-to-token circuit en tendencias como copying/anti-copying, y revela que el modelo de una capa depende fuertemente de un copying extendido. Depende de que el mapeo sea sobre el mismo espacio y de una estructura aproximadamente lineal; frente al MLP y a modelos grandes, solo sirve como un indicio limitado.
